%% ============================================================================
%% III. Problem Statement                                     target: ~0.5 page
%% Stage 1: the problem statement must be fully developed (a concrete technical
%% formulation, not a high-level argument).
%% ============================================================================
\section{Preliminaries}
\label{sec:Preliminaries}

\begin{figure}[t]
  \centering
  \figplaceholder{\linewidth}{3.6cm}{%
    \textbf{Fig.~2 (notation).} Gripper pose $\gpose$, rod of length $L$,
    contact at arc length $s_{\mathrm{c}}$ and ring angle $\cang$, wrist wrench
    $\wrench$, tip pose $\tpose$, and target $\tposegoal$.}
  \caption{Problem setup and notation.}
  \label{fig:notation}
\end{figure}

\begin{outline}
  \item \textbf{Object.} A naturally straight, superelastic Nitinol rod of
    length $L$, modeled as an inextensible Kirchhoff rod under quasi-static
    manipulation~\cite{bretl2014quasistatic}, grasped rigidly at one end; the
    free tip carries a tool. \todo{Rod length, diameter, material constants.}
  \item \textbf{Robot and sensing.} Gripper pose $\gpose\in\SE{3}$ ($\SE{2}$
    in planar experiments) from forward kinematics; wrist wrench
    $\wrench=(\vect{f},\vect{\tau})\in\R^{6}$ from a force/torque sensor; a
    sensorized environment edge (ring) that reports contact and the contact
    angle $\cang$ along the ring (\cref{fig:notation}).
  \item \textbf{Contact.} The rod touches the edge at a single point at an
    unknown arc length $s_{\mathrm{c}}\in(0,L)$; the edge applies an unknown
    reaction force (normal and friction). \todo{Sticking vs.\ sliding contact;
    friction model.}
  \item \textbf{Observation} at time $t$: $\vect{o}_t=(\gpose,\wrench,\cpt)$,
    where $\cpt$ holds the contact state and $\cang$. No vision is used during
    execution; an RGB-D camera provides ground truth only.
  \item \textbf{Reachability.} Let $\Rfree$ and $\Rcon$ be the sets of tip
    poses reachable in the workspace $\mathcal{W}$ without and with an edge
    contact, respectively.
  \item \textbf{Assumptions.} Quasi-static motion, a single contact, known edge
    geometry, deformation within the rod's elastic range.
    \todo{Gravity: modeled or negligible?}
\end{outline}

\begin{problem}[Haptic tip-pose estimation]\label{prob:estimation}
Given the observation $\vect{o}_t$, estimate the tip pose, $\tposehat\approx\tpose$.
\end{problem}

\begin{problem}[Contact-aided tip manipulation]\label{prob:manipulation}
Given a target $\tposegoal\in\Rcon\setminus\Rfree$, find gripper motions that
bring $\tpose$ to within position and orientation tolerances
$(\epsilon_p,\epsilon_\theta)$ of $\tposegoal$ by using a contact with the
edge, while avoiding other collisions and keeping the rod within its elastic
range.
\end{problem}
